% !TeX root = ../papers/nc-p3.tex
\section{滑らかな種数二の階数一五次曲線の排除}
\label{sec:rank-one-genus-two-quintic}
\begin{theorem}\label{thm:rank-one-genus-two-quintic}
被約な四辺有限の標準有限Frobenius窓の一般可逆な全階数反対根において、
階数一軌跡は滑らかな種数二の整五次曲線を含まない。
\end{theorem}
\begin{proof}[二つの完全五次線型系]
そのような曲線を $C$ と仮定する。前定理の総次数から他の曲線成分はなく、二つの埋込みの像は非平面的である。
超平面束を $L,L'=L\eta$ と書けば $\eta^2=\mathcal O$ である。
次数五、種数二からどちらの一次系も四切断の完全系である。二倍束の切断数は9で、二つの独立な二次式は
非平面五次曲線を含めないから、各像は唯一の既約二次曲面 $Q,Q'$ に載る。
$\eta=\mathcal O$ なら二つの埋込みは射影同値であり、線型座標を合わせると $C$ 上で $N=xx^t$ となる
（同じ線束の基底同一視と定数倍を取る）。従って $N=xx^t+qA$ で $q$ は $Q$ の式、$A$ は定数対称行列である。
$q$ の外の印により $A$ の階数は2以下となり一般可逆性に反する。
以後 $\eta$ は非自明である。$C$ の外の指定印数を $s$ とする。
十六の双一次式は全印上独立、$h^0(C,LL')=9$ なので $s\ge7$。

以下、平面内の出力点の平方の線型関係を入力で読む。
幾何学的な異なる $m\le7$ 点の関係で係数の非零な台を取ると、同じ台の入力一次座標空間は次元
$\lfloor m/2\rfloor$ 以下である。二次Veronese間の線型同型により同じ関係が非零係数の再規格化を除き入力の平方にも成立し、
一次座標の評価像が非退化対角対合で全等方になるからである。一般の平面族の関数体上で関係を取れば、
分離的な一つの拡大体から来る動点について、係数は全部零か全て非零になる（共役を取る）。
\end{proof}
\begin{proof}[孤立印と平面鉛筆]
曲線外の出力印 $p'$ は $Q'$ に載らない。もし載ればその点を通る一般平面が出力曲線を五点で横断し、
これらと $p'$ は同じ二次曲線に入るので六平方には関係がある。五つの動点は一つの体の共役である。
実際平面入射は曲線上の射影直線束で既約、一般平面への次数は5である。関係は $p'$ のみには支えを持てず、
従って対応する入力五点は共面となる。この五点は出力超平面の完全な切断因子で、入力でも次数5なので
$L'=L$ に反する。曲線外の印が全て $Q'$ の外にあることを得た。

二つの曲線外の印を出力 $p',q'$、入力 $p,q$ とし、出力割線 $J'$ と曲線の交わりの長さを $b$ とする。
$p'$ は $Q'$ 外なので $b\le2$ である。$J'$ を通る平面鉛筆の一般元を $H'$ とし、
$D=C\cap H'$ を出力で切った長さ5の因子とする。平面 $H'$ の二次式で $D$ を含むものは $Q'|_{H'}$ の倍のみである。
実際線型正規性より $H^1(\mathcal I_C(1))=0$ なので超平面切断完全列で全ての二次式が持ち上がる。
従って $D$ の五平方は階数5で、$p'$ または $q'$ を一つ加えると階数6になる。

$b=0,1$ の場合、一般 $D$ は被約である。固定部分は長さ $b$ で残りは次数 $5-b$ の射 $C\to\mathbb P^1$ の一般繊維であり、
標数零の一つの分離拡大の共役点である。$D,p',q'$ の七平方の関係は $p',q'$ の両方で非零、
かつ動点でも全て非零である。固定部分と $p',q'$ だけでは直線上の高々三点で二次評価が独立だからである。
入力では動点と $p,q$ が一平面に入る。$b=0$ なら再び五点の因子から $L=L'$ となり不可能である。

$b=1$ として固定点を $r\in C$ とする。入力割線 $\langle p,q\rangle$ からの曲線の射影は、
出力鉛筆の一般四点繊維で一定である。それらと $p,q$ が共面で、一般動点は入力割線に載らないためである。
両射は非定数で、分離拡大の共役値が全て一致するから入力鉛筆の関数は出力鉛筆の関数体に属する。
入力射影の次数は4の倍数で5以下、従って4で、二つの鉛筆は底の射影変更で一致する。
入力鉛筆の固定因子は一点 $z$、ゆえに $L(-z)=L'(-r)$ となる。
非自明な二捩れなので $z\ne r$ かつ $2z\sim2r$ である。これは種数二の唯一の次数二鉛筆
（Riemann--Rochより二切断を持つ次数二線束は標準束）での二つのWeierstrass点である。
固定した $\eta$ にこのような順序を忘れた二点組は高々一つである。
実際 $z-r\sim z'-r'$ なら $z+r'\sim z'+r$ であり、二因子が等しくなければ両方とも標準鉛筆の繊維である。
Weierstrass点を含む繊維はその二重点なので、この場合にも同一の二点組となる。
\end{proof}
\begin{proof}[各曲線外の印から見た割線の個数]
$b=1$ の割線に三つの曲線外の印があることはない。あれば $r$ と合わせた四つの相異なる共線出力点には
二次式で全係数非零の関係があり、上の全等方性から入力四点も共線である。すると入力割線の固定点 $z$ が
$r$ になり矛盾する。従って各 $p'$ から $b=1$ の割線が結べる他の印は高々二つで、各Weierstrass点を通る方向に一つずつである。

$b=2$ の割線にも曲線外の印は高々二つである。三つあれば出力直線上で二つの曲線交点（接する長さ二も含む）と
三印の和が長さ5になる。逆向きの二次Veronese座標変換に対する二行小行列式は次数4でその全てで消える。
よって出力直線全体が階数一曲線になり、既に五次曲線に使った総次数を超える。
また固定した $p'$ からの $b=2$ の方向は四つ以下である。
$p'$ からの平面射影は次数5の超平面束による基点のない射で、像は非退化なので5の素数性より平面五次への双有理射となる。
滑らかな $C$ はその正規化であり、特異点のデルタ不変量の総和は $6-2=4$。
長さ二の割線方向は二枝または非浸入枝の像であり、それぞれ少なくとも1を使う。

各 $p'$ に対し $s-1\le 2+4$ を得る。既に $s\ge7$ なので $s=7$ かつ全て等号である。
従って二つのWeierstrass点の一方を固定すると、七つの印はそれぞれその点を通る $b=1$ の割線上に
ただ一つ別の印を持つ。これは七元集合を二元集合に分割してしまう。以上で種数二を排除した。
\end{proof}
